\begin{landscape}
\begin{table}[!h]
\begin{center}
\bgroup
\def\arraystretch{1.5}
    \begin{tabular}{| p{3cm} | p{5cm} | p{6cm} | p{2cm} | p{6cm} |}
    \hline
\textbf{Stakeholder} & \textbf{Hauptinteressen am Projekt} & \textbf{Erwartungen / Anforderungen (erste Ideen)} & \textbf{Einfluss} & \textbf{Potenzielle Konflikte} \\
\hline
\hline
\textbf{Bürger:in} & Einfache Bedienung, schnelles Feedback, erfolgreiche Platzierung & Klare Anweisungen. verständliche Fehlermeldungen & Mittel & Wunsch nach maximaler Freiheit vs. restriktive Regeln \\
\hline
\textbf{Stadtplaner:in} & Einhaltung städtischer Vorgaben, weniger ungültige Vorschläge & Flexibles Regelmanagement, präzise Validierung & Hoch & Komplexität der Regeln vs. Performance / Verständlichkeit für Bürger:innen \\
\hline
\textbf{Politische Entscheidungsträger:in} & Sichtbare Bürgerbeteiligung, positives Image, rechtliche Absicherung & DSGVO-Konformität, möglichst geringer Pflegeaufwand & Hoch & Offene Beteiligung vs. Kontrollverlust / Budget- und Zeitrahmen \\
\hline
\textbf{Entwickler:in} & Stabile Technik, geringe Latenz, überschaubarer Integrationsaufwand & Typsichere API, klare Fehlercodes, Versionierung & Mittel & Komplexe Regelsets vs. Performance-Budget \\
\hline
\textbf{Feuerwehr und Infrastruktur-Anbieter} & Vorschläge sollten nicht in operativen Betrieb eingreifen (Wendekreis, Wartung, …) & Modellierung von Nutzungsbereichen & Mittel & Einfachheit der Wartung vs. Freiheit von Vorschlägen \\
\hline
\end{tabular}
    \egroup
    
\caption{Stakeholder-Tabelle}
\label{tab:Stakeholder-Tabelle}
\end{center}
\end{table}
\end{landscape}
\FloatBarrier